\begin{definition}[벡터의 합과 스칼라 곱의 성질]\label{def:vector1}
    \begin{itemize}
        \item 모든 벡터 $\mathbf{x}$, $\mathbf{y}$에 대해서 $\mathbf{x} + \mathbf{y} = \mathbf{y} + \mathbf{x}$이다.
        \item 모든 벡터 $\mathbf{x, y
        , z}$에 대해서 $(\mathbf{x} + \mathbf{y}) + \mathbf{z} = \mathbf{x} + (\mathbf{y} + \mathbf{z})$이다. 
        \item 모든 벡터 $\mathbf{x}$에 대해서 $\mathbf{x} + 0 = \mathbf{x}$를 만족하는 벡터 0이 존재한다. 
        \item 각 벡터 $\mathbf{x}$마다 $\mathbf{x} + \mathbf{y} = 0$을 만족하는 벡터 $\mathbf{y}$가 존재한다. 
        \item 모든 벡터 $\mathbf{x}$에 대하여 $1 \mathbf{x} = \mathbf{x}$이다.
        \item 모든 실수 $a, b$와 모든 벡터 $\mathbf{x}$에 대해서 $(\mathbf{ab})\mathbf{x} = \mathbf{a}(\mathbf{bx})$이다. 
        \item 모든 실수 $a$와 모든 벡터 $\mathbf{x, y}$에 대해서 $a(\mathbf{x}+ \mathbf{y}) = a \mathbf{x} + a \mathbf{y}$이다. 
        \item 모든 실수 $a, b$와 모든 벡터 $\mathbf{x}$에 대하여 $(a+b) \mathbf{x} = a \mathbf{x} + b \mathbf{x}$이다.
    \end{itemize}\footnote{출처 : 프리드버그의 선형대수학}
\end{definition}

\defn{벡터의 합과 스칼라 곱의 성질\label{def:vector1}}{\begin{itemize}
        \item 모든 벡터 $\mathbf{x}$, $\mathbf{y}$에 대해서 $\mathbf{x} + \mathbf{y} = \mathbf{y} + \mathbf{x}$이다.
        \item 모든 벡터 $\mathbf{x, y, z}$에 대해서 $(\mathbf{x} + \mathbf{y}) + \mathbf{z} = \mathbf{x} + (\mathbf{y} + \mathbf{z})$이다. 
        \item 모든 벡터 $\mathbf{x}$에 대해서 $\mathbf{x} + 0 = \mathbf{x}$를 만족하는 벡터 0이 존재한다. 
        \item 각 벡터 $\mathbf{x}$마다 $\mathbf{x} + \mathbf{y} = 0$을 만족하는 벡터 $\mathbf{y}$가 존재한다. 
        \item 모든 벡터 $\mathbf{x}$에 대하여 $1 \mathbf{x} = \mathbf{x}$이다.
        \item 모든 실수 $a, b$와 모든 벡터 $\mathbf{x}$에 대해서 $(\mathbf{ab})\mathbf{x} = \mathbf{a}(\mathbf{bx})$이다. 
        \item 모든 실수 $a$와 모든 벡터 $\mathbf{x, y}$에 대해서 $a(\mathbf{x}+ \mathbf{y}) = a \mathbf{x} + a \mathbf{y}$이다. 
        \item 모든 실수 $a, b$와 모든 벡터 $\mathbf{x}$에 대하여 $(a+b) \mathbf{x} = a \mathbf{x} + b \mathbf{x}$이다.
    \end{itemize}\footnote{출처 : 프리드버그의 선형대수학}}


    ; find \begin{definition} as $all >> \defn{$1\label{$2}}{$3} where $all matches
/\\begin\{definition\}\[(.*?)\]\s*\\label\{(.*?)\}\s*([\s\S]*?)\\end\{definition\}/